\section{Exact Chinese-Remainder Reconstruction}
\label{sec:crt}

The inverse problem is local.  Given the two child residues of a parent class,
one must reconstruct the unique parent residue without solving a global linear
system.  The required Bezout coefficient has a closed form for every Bruun
sibling pair.

\subsection{Trinomial siblings}

Let $h$ be even, $b\ne0$, and
\begin{equation}
 p_-=z^h-bz^{h/2}+1,
 \qquad
 p_+=z^h+bz^{h/2}+1.
 \label{eq:crt-siblings}
\end{equation}

\begin{lemma}[Local Bezout coefficient]
In $\K_N[z]/(p_+)$,
\begin{equation}
 p_-^{-1}\equiv {z^{h/2}+b\over 2b}.
 \label{eq:local-bezout}
\end{equation}
\end{lemma}
\begin{proof}
Modulo $p_+$, multiplication by $z^{-h/2}$ gives
$z^{-h/2}\equiv-(z^{h/2}+b)$.  Also
$p_--p_+=-2bz^{h/2}$.  Hence
$p_-(z^{h/2}+b)/(2b)\equiv1\pmod{p_+}$.
\end{proof}

\begin{proposition}[Explicit sibling merge]
Let $r_-$ and $r_+$ be residues modulo $p_-$ and $p_+$, respectively.  Put
\begin{equation}
 t=\left[
 { (r_+-r_-)(z^{h/2}+b)\over 2b}
 \right]_{p_+},
 \qquad
 c=[r_-+p_-t]_{p_-p_+}.
 \label{eq:trinomial-merge}
\end{equation}
Then $c\equiv r_-\pmod{p_-}$ and $c\equiv r_+\pmod{p_+}$, and $c$ is unique
modulo $p_-p_+$.
\end{proposition}
\begin{proof}
Choose the canonical polynomial representatives of degrees below $h$ in the
two child quotients.  The expression in \eqref{eq:trinomial-merge} then defines
a polynomial lift; changing either lift by its child factor changes $c$ only by
a multiple of $p_-p_+$.
The first congruence is immediate.  For the second, use
\eqref{eq:local-bezout} to obtain
$p_-t\equiv r_+-r_-\pmod{p_+}$.  Coprimality and the Chinese remainder theorem
give uniqueness.
\end{proof}

At a binomial node, $p_-=z^h-1$ and $p_+=z^h+1$.  The corresponding merge is
\begin{equation}
 t={r_--r_+\over2},
 \qquad c=r_-+(z^h-1)t.
 \label{eq:binomial-merge}
\end{equation}
The factors $1/2$ and $1/(2b)$ are the local CRT interpolation measures.  Their
relation to the uniform inverse Fourier scale is established only after the
normalized packed map is identified in Section~\ref{sec:fourier-identity}.

\subsection{Global inverse and fixed points}

Let $A_N$ denote the product algebra on the right side of
\eqref{eq:global-reduction}.  Define $\rec_N:A_N\to R_N$ by applying
\eqref{eq:trinomial-merge} and \eqref{eq:binomial-merge} from leaves to root.

\begin{theorem}[Two-sided Bruun reconstruction]
The reduction and reconstruction maps are mutual inverses:
\begin{equation}
 \rec_N\red_N=\id_{R_N},
 \qquad
 \red_N\rec_N=\id_{A_N}.
 \label{eq:two-sided-crt}
\end{equation}
\end{theorem}
\begin{proof}
At one node, the explicit merge is the inverse of the residue map by the
preceding proposition and the binomial identity.  Induction on the height of
the factor tree proves both compositions in \eqref{eq:two-sided-crt}.
\end{proof}

\begin{corollary}[Exact fixed-point law]
Every polynomial state is a fixed point of reconstruction after reduction, and
every admissible leaf state is a fixed point of reduction after reconstruction.
These are identities over $\R$; their structure constants lie in $\K_N$.
They are not convergence statements.
\end{corollary}

\begin{example}[A symbolic $N=8$ merge]
The siblings $p_-=z^2-\sqrt2z+1$ and
$p_+=z^2+\sqrt2z+1$ have $h=2$ and $b=\sqrt2$.  For linear residues
$r_-=a_0+a_1z$ and $r_+=b_0+b_1z$, the unique quartic parent is
\begin{equation}
 c=r_-+p_-\left[
 {(r_+-r_-)(z+\sqrt2)\over2\sqrt2}
 \right]_{p_+}.
 \label{eq:n8-symbolic-merge}
\end{equation}
Substitution modulo either child recovers its prescribed residue.  Applying the
binomial merge above this node recovers the degree-seven input class.
\end{example}
